% Same code as the original growth problem; qualitative solutions for Fall 2026.
\question \textbf{Fast Exponentiation:} in this problem, we will examine a
real-world algorithm used to improve the speed of calculating exponents.
You can assume for the purposes of this problem that multiplication is a
fixed-cost operation. Classify each version as linear, sublinear, or
superlinear relative to the non-negative integer exponent \lstinline{n};
explain your answer in words.

\begin{parts}

\begin{blocksection}
\part First, determine the runtime efficiency of the naive exponentiation algorithm.
\begin{lstlisting}
def exp(b, n):
    if n == 0:
        return 1
    else:
        return b * exp(b, n - 1)
\end{lstlisting}
\begin{solution}[0.25in]
Linear. Each call reduces \lstinline{n} by one. Doubling the exponent
roughly doubles the number of calls and multiplications.
\end{solution}
\end{blocksection}

\begin{blocksection}
\part Now, express the runtime of the fast exponentiation algorithm.
\begin{lstlisting}
def fast_exp(b, n):
    if n == 0:
        return 1
    elif n % 2 == 0: # Assume square runs in constant time
        return square(fast_exp(b, n // 2))
    else:
        return b * fast_exp(b, n - 1)
\end{lstlisting}
\begin{solution}[0.25in]
Sublinear. An even exponent is halved; an odd exponent is reduced by one
and then halved on the next step. Doubling the starting exponent adds only
one initial halving step. The recursive result is computed once and then
passed to \lstinline{square}.
\end{solution}
\end{blocksection}

\begin{blocksection}
\part What about this slightly modified version of \texttt{fast\char`_exp}?
\begin{lstlisting}
def fast_exp(b, n):
    for _ in range(50 * n):
        print("Killing time")
    if n == 0:
        return 1
    elif n % 2 == 0:
        return square(fast_exp(b, n // 2))
    else:
        return b * fast_exp(b, n - 1)
\end{lstlisting}
\begin{solution}[0.25in]
Linear. The first loop alone does work proportional to \lstinline{n}.
Subsequent loops shrink as the exponent is halved. Odd exponents add one
extra stage before a halving, but the total work still stays proportional
to the starting exponent. The fixed multiplier 50 does not change the
category.
\end{solution}
\end{blocksection}

\end{parts}

\begin{guide}
\textbf{Teaching Tips}
\begin{itemize}
    \item Compare small exponents and describe what changes when the exponent
    doubles. No formal bounds or recurrence equations are needed.
    \item Distinguish reusing a local recursive result from keeping a cache
    across calls. Fast exponentiation reduces work by halving the problem;
    it does not itself maintain a persistent cache.
    \item These functions do not repeat subproblems within a call, so a cache
    mainly helps when the same inputs are requested again. A cache for a
    two-argument function must distinguish both the base and the exponent.
\end{itemize}
\end{guide}
